\documentclass[../differential_methods.tex]{subfiles}
\begin{document}
\section{Aula 21 - 22 de Junho, 2026}
\subsection{Motivações}
\begin{itemize}
	\item Equações da onda.
\end{itemize}
\subsection{Equação da Onda}
Consideremos a equação da onda
\[
	u_{tt} - c^{2}u_{xx} = 0, \quad x\in [0, 1],
\]
com condições iniciais
\[
	u(x, 0) = f(x),\quad  u_{t}(x, 0) = g(x)
\]
e fronteiras
\[
	u(0, t) = h_1(t), \quad u(1, t)=h_2(t).
\]
Vamos definir \(v = cu_x\) e \(w=-u_t\); assim, obtemos as equações
\[
	v_{t} = cu_{tt},\; v_{x} = cu_{xx},\; w_{x} = -u_{tt} \;\&\; w_{t} = -u_{tt},
\]
ou seja, obtivemos o sistema acoplado
\[
	\left\{\begin{array}{ll}
		v_t + cw_x =0                 \\
		w_t + cv_x = 0                \\
		v(x, 0) = cu_x(x, 0) = cf'(x) \\
		w(x, 0) = -u_t(x, 0) = -g(x)
	\end{array}\right.
\]
que, em forma de matriz, assume
\[
	\frac{\partial^{}}{\partial t^{}}\begin{bmatrix}
		v \\
		w
	\end{bmatrix} + \frac{\partial^{}}{\partial x^{}} \begin{bmatrix}
		0 & c \\
		c & 0
	\end{bmatrix} \begin{bmatrix}
		v \\
		w
	\end{bmatrix} = 0, \quad \underline{u}(x, 0) = \begin{bmatrix}
		cf'(x) \\
		-g(x)
	\end{bmatrix},
\]
e chegamos na equação
\[
	\underline{u}_t + A \underline{u}_x = 0.
\]
Agora, diagonalizamos A por meio dos seus autovalores e autovetores para desacoplar o sistema:
\[
	\det{(A-\lambda I)}=0 \Rightarrow \det \begin{bmatrix}
		-\lambda & c
		c        & -\lambda
	\end{bmatrix} =0 \Rightarrow \lambda^{2} - c^{2}=0 \Rightarrow \lambda_1 = -c, \; \lambda_2 = c.
\]
Os autovetores, por sua vez, são
\[
	\left\{\begin{array}{ll}
		A \underline{v}^{1} = -c \underline{v}^{1} \\
		A \underline{v}^{2} = -c \underline{v}^{2}
	\end{array}\right. \Rightarrow \underline{v}^{1}=\begin{bmatrix}
		-1 \\
		1
	\end{bmatrix} \quad\&\quad \underline{v}^{2} = \begin{bmatrix}
		1 \\
		1
	\end{bmatrix}
\]
Sendo assim, colocando
\[
	R = \begin{bmatrix}
		-1 & 1 \\
		1  & 1
	\end{bmatrix} \Rightarrow R^{-1}=\begin{bmatrix}
		-\frac{1}{2} & \frac{1}{2} \\
		\frac{1}{2}  & \frac{1}{2}
	\end{bmatrix},\; \Lambda = \begin{bmatrix}
		-c & 0 \\
		0  & c
	\end{bmatrix},
\]
obtemos a diagonalização de A como \(A = R\Lambda R^{-1}\). Na equação diferencial, obtemos
\[
	\underline{u}_t + A \underline{u}_x = 0 \Rightarrow \underline{u}_t + R\Lambda R^{-1}\underline{u}_x = 0 \Rightarrow R^{-1}\underline{u}_t + \Lambda R^{-1}\underline{u}_x = 0,
\]
tal que, colocando \(\underline{w} \coloneqq R^{-1}\underline{u}\) com coordenadas \(w_1\) e \(w_2\), a equação torna-se
\[
	\underline{w}_t + \Lambda \underline{w}_x = 0 \Rightarrow \left\{\begin{array}{ll}
		\frac{\partial^{}}{\partial t^{}}w_1 - c \frac{\partial^{}}{\partial x^{}}w_1 = 0 \\
		\frac{\partial^{}}{\partial t^{}}w_2 + c \frac{\partial^{}}{\partial x^{}}w_2 = 0
	\end{array}\right.
\]
Pra que todo esse esforço? Pois estes \(w_1\) e \(w_2\) desacoplados podem ser vistos como ondas fundamentais da solução, e são duas equações de transporte com sinais opostos: a primeira transporta a onda \(w_1\) para a esquerda, e a outra transporta a \(w_2\) para a direita -- a solução da onda nada mais é do que duas soluções da equação do transporte!
Uma delas gera características indo pra esquerda, a outra gera pra direita, sem que haja sobreposição de informações!

Construir a solução exata dessa equação se dá por um processo semelhante ao das equações dos transportes: considerando cada uma isolada, sabemos que a solução se dá pela condição inicial, ou seja,
\[
	w_1(x, t) = w_1(x+ct, 0) \quad\&\quad w_2(x, t) = w_2(x- ct, 0),
\]
mas isso requer que as condições iniciais do sistema que usamos no começo também passem pelas transformações aplicadas (depois de resolver, pode reverter as transformações pelas inversas).

Finalmente, podemos estimar a solução numérica dos sistemas acoplado (Sistema 1) e desacoplado (Sistema 2); este segundo, por ser só duas equações do transporte, já temos uma boa noção do que fazer -- por exemplo, podemos aplicar Lax-Friedrischs, utilizando as notações
\[
	w_1(x_{j}, t_{n})\cong W_{1}\biggl|_{j}^{n}\biggr. \quad\& \quad w_2(x_{j}, t_{n})\cong W_{2}\biggl|_{j}^{n}\biggr.
\]
para obter
\[
	W_{1}\biggl|_{j}^{n+1}\biggr. = \frac{1}{2}(W_{1} \biggl|_{j+1}^{n}\biggr. + W_1 \biggl|_{j-1}^{n}\biggr.) - (-c)\frac{\Delta t}{2\Delta x}(W_1 \biggl|_{j+1}^{n}\biggr. - W_1 \biggl|_{j-1}^{n}\biggr.)
\]
e
\[
	W_{2}\biggl|_{j}^{n+1}\biggr. = \frac{1}{2}(W_{2} \biggl|_{j+1}^{n}\biggr. + W_2 \biggl|_{j-1}^{n}\biggr.) - (-c)\frac{\Delta t}{2\Delta x}(W_2 \biggl|_{j+1}^{n}\biggr. - W_2 \biggl|_{j-1}^{n}\biggr.),
\]
com forma matricial
\[
	\begin{bmatrix}
		W_1 \bigl|_{j}^{n+1}\bigr. \\
		W_2 \bigl|_{j}^{n+1}\bigr.
	\end{bmatrix} = \frac{1}{2}\biggl(\begin{bmatrix}
			W_1 \bigl|_{j+1}^{n}\bigr. \\
			W_2 \bigl|_{j+1}^{n}\bigr.
		\end{bmatrix}+\begin{bmatrix}
			W_1 \bigl|_{j-1}^{n}\bigr. \\
			W_2 \bigl|_{j-1}^{n}\bigr.
		\end{bmatrix} \biggr) - \Lambda \frac{\Delta t}{2\Delta x}\biggl(\begin{bmatrix}
			W_1 \bigl|_{j+1}^{n}\bigr. \\
			W_2 \bigl|_{j+1}^{n}\bigr.
		\end{bmatrix} - \begin{bmatrix}
			W_1 \bigl|_{j-1}^{n}\bigr. \\
			W_2 \bigl|_{j-1}^{n}\bigr.
		\end{bmatrix}\biggr),
\]
ou, equivalentemente,
\[
	\underline{W}_{j}^{n+1} = \frac{1}{2}(\underline{W}_{j+1}^{n}+\underline{W}_{j-1}^{n})-\Lambda \frac{\Delta t}{2\Delta x}(\underline{W}_{j+1}^{n}-\underline{W}_{j-1}^{n}),
\]
que podemos multiplicar por R para chegar em
\begin{align*}
	                    & R\underline{W}_{j}^{n+1} = \frac{1}{2}R(\underline{W}_{j+1}^{n}+\underline{W}_{j-1}^{n})-R\Lambda \frac{\Delta t}{2\Delta x}(\underline{W}_{j+1}^{n}-\underline{W}_{j-1}^{n})                   \\
	\Longleftrightarrow & RR^{-1}\underline{W}_{j}^{n+1} = \frac{1}{2}RR^{-1}(\underline{W}_{j+1}^{n}+\underline{W}_{j-1}^{n})-RR^{-1}\Lambda \frac{\Delta t}{2\Delta x}(\underline{W}_{j+1}^{n}-\underline{W}_{j-1}^{n}) \\
	\Longleftrightarrow & \boxed{\underline{U}_{j}^{n+1} = \frac{1}{2}(U_{j+1}^{n}+U_{j-1}^{n}) - A \frac{\Delta t}{2\Delta x}(\underline{U}_{j+1}^{n} - \underline{U}_{j-1}^{n})}
\end{align*}
Com \(\underline{U}_{j}^{n} = \begin{bmatrix}
	V_{j}^{n} \\
	W_{j}^{n}
\end{bmatrix}\), recuperamos o \hyperlink{lax_friedrichs}{\textit{Lax-Friedrichs}} de antes, com a única diferença de que o escalar ``a'' multiplicando tornou-se uma matriz ``A''! Similarmente, o \hyperlink{lax_wandroff}{\textit{Lax-Wandroff}} assume a forma
\[
	\boxed{\underline{U}_{j}^{n+1} = \underline{U}_{j}^{n} - \frac{\Delta t}{2\Delta x} A (\underline{U}_{j+1}^{n} - \underline{U}_{j-1}^{n}) + \frac{\Delta t^{2}}{2\Delta x^{2}} A^{2}(\underline{U}_{j+1}^{n}-2\underline{U}_{j}^{n}+\underline{U}_{j-1}^{n})}
\]
e o \hyperlink{up_wind}{\textit{upwind}} também, mas com uma pequena mudança de que a aproximação pra \(w_1\) e \(w_2\) têm sinais opostos, então separamos a matriz em uma só para a parte positiva e outra só para a parte negativa:
\[
	\Lambda  = \begin{bmatrix}
		-c & 0 \\
		0  & c
	\end{bmatrix} = \begin{bmatrix}
		0 & 0 \\
		0 & c
	\end{bmatrix} + \begin{bmatrix}
		-c & 0 \\
		0  & 0
	\end{bmatrix} \eqqcolon \Lambda^{+} + \Lambda^{-},
\]
fornecendo a aproximação para \(\underline{W}_{j}^{n+1}\) na forma
\[
	\underline{W}_{j}^{n+1} = \underline{W}_{i}^{n} - \frac{\Delta t}{\Delta x}\Lambda^{+}(\underline{W}_{j}^{n} - \underline{W}_{j-1}^{n}) - \frac{\Delta t}{\Delta x}\Lambda^{-}(\underline{W}_{j+1}^{n} - \underline{W}_{j}^{n}).
\]
\begin{exr}
	Faça a transformação que fizemos previamente para obter a forma do sistema acoplado, com \(A^{+} \coloneqq R\Lambda^{+}R^{-1}\) e \(A^{-}\coloneqq R\Lambda ^{-}R^{-1}\).
\end{exr}

Existe um método mais direto para resolver, partindo da mesma equação definida no começo,
\[
	u_{tt} - c^{2}u_{xx} = 0, \quad x\in [0, 1],
\]
com condições iniciais
\[
	u(x, 0) = f(x),\quad  u_{t}(x, 0) = g(x)
\]
e contorno
\[
	u(0, t) = h_1(t), \quad u(1, t)=h_2(t).
\]
Dessa vez, faremos a estratégia de discretizar em x e em t com 2 passos e níveis (a discretização em cruz), com todo aquele processo que já conhecemos de avaliar em \((x_{j}, t_{n})\), substituir a derivada, descartar o erro de truncamento, etc, resultando em
\[
	\frac{1}{\Delta t^{2}}(U_{j}^{n+1} - 2 U_{j}^{n} + U_{j}^{n-1}) - c^{2} \frac{1}{\Delta x^{2}}(U_{j+1}^{n}-2U_{j}^{n} + U_{j-1}^{n}) = 0.
\]
Reescrevendo ela, com \(\sigma^{2} \coloneqq c^{2}\frac{\Delta t^{2}}{\Delta x^{2}}\), chegamos em
\[
	U_{j}^{n+1} = -U_{j}^{n-1} + \sigma^{2}U_{j-1}^{n} + (2-2\sigma ^{2})U_{j}^{n} + \sigma^{2}U_{j+1}^{n},
\]
que normalmente cairia na questão de ter um problema no limite de tempo mais no passado ou mais no futuro, mas podemos nos aproveitar da condição \(u_t(x, 0) = g(x)\) para lidar com isso: quando \(n=0\), a equação seria
\[
	U_{j}^{1} = - U_{j}^{-1} + \sigma^{2} U_{j-1}^{0} + (2-2\sigma^{2})U_{j}^{0} + \sigma^{2}U_{j+1}^{0};
\]
todos os termos \(U^{0}\) são da forma \(f(x_{j})\) pela condição inicial, enquanto que o \(-U_{j}^{-1}\) pode ser lidado com utilizando a derivada:
\[
	U_t(x_{j}, 0) = g(x_{j}) \Rightarrow \frac{U_{j}^{1} - U_{j}^{-1}}{2\Delta t} = g(x_{j}) \Rightarrow U_{j}^{-1} = U_{j}^{1} - 2\Delta tg(x_{j}),
\]
a partir do qual basta substituir na equação original:
\begin{align*}
	            & U_{j}^{1} = - U_{j}^{1} + 2\Delta tg(x_{j}) + \sigma^{2} U_{j-1}^{0} + (2-2\sigma^{2})U_{j}^{0} + \sigma^{2}U_{j+1}^{0} \\
	\Rightarrow & \boxed{2U_{j}^{1} = 2\Delta tg(x_{j}) + \sigma^{2} U_{j-1}^{0} + (2-2\sigma^{2})U_{j}^{0} + \sigma^{2} U_{j+1}^{0}}
\end{align*}
Agora, fica a grande pergunta -- este método funciona? Descubra!
\begin{exr}
	Analise a convergência do método: compute o erro de truncamento, mostre que a consistência dele é da ordem \(\mathcal{O}(\Delta t^{2}+\Delta x^{2})\) e que a estabilidade dele por Von Neumann é dada pela condição \(| c |\frac{\Delta t}{\Delta x}\leq 1\)
\end{exr}
Surpreendentemente, o método de diferenças centrais funciona bem aqui, mesmo sendo horrível para a equação do transporte! Moral da história: se for propor um método diferente, tem que analisar ele pra ver se vai funcionar ou não, independente do que você ache que vá acontecer.

Um último comentário é da questão da condição CFL, que mantém-se na forma \(\frac{| c |\Delta t}{\Delta x} \leq 1\), e a condição é satisfeita para que haja a convergência.

\end{document}
